% s-172-p172-m002; original PDF 172, 173, 174
\Aphorism{s-172-p172-m002}{259}
\Sperr{Trans Caucasum.} Man macht das Gesicht zu seinen Gedanken, oder — man hat die Gedanken zu seinem Gesicht. Die feineren Beobachter entscheiden sich für die zweite Lesart. Wir Europäer setzen unseren Stolz darein, vielerlei Gesichter von oft nur differentieller Abweichung zu bringen. Demgemäss befähigen wir uns auch, vielerlei Meinungen zu haben und oft sehr feine, nüancirte Meinungen, die nicht in die ersten besten Worte einzufangen sind. Daneben haben wir unsere Standpunkte, die eigentlich keine Meinungen, sondern Mittelwerthe von Meinungen sind; starke Gedankenbündel, mit denen sich besser schlagen lässt als mit den einzelnen Ruthen. Auch zu den Standpunkten werden wir physiognomisch ermächtigt, und der bitterste Kampf der Überzeugungen ist nur eine Oberflächenerscheinung zu dem darunter glimmenden Antagonismus der Gesichter. Eine industrielle Nase und zwei Schwärmeraugen können einander nicht verstehen; und der Engländer, der sein Conterfei in die Tiefe denkt, ist in alle Ewigkeit vor deutscher Philosophie behütet. So ist Europa ein Kampf von Angesicht zu Angesicht, von Meinung zu Meinung — so scheint es. Und doch, wieviel Gleichartigkeit noch in diesem Wirrwarr, wieviel Axiom in dieser „freien Forschung“, welche Harmonie in diesem sich befehdenden und missverstehenden Europa! Ein Europäer mag Germane oder Lateiner sein, Christ oder Jude, Phantast oder Realist: er mag sich stellen und werfen wie er will — zum Schluss fällt und steht er doch auf seinem Europa, auf ein, zwei, drei Grundwahrheiten, die zwischen den Meridianen von Madrid und Moskau unumstösslich geglaubt werden: ich werde mich hüten, diese Grundwahrheiten noch einmal zu sagen. Woran liegt es, wer garantirt diese Einheit im Glaubenskern, bei aller Sektirer-Verschiedenheit in der Schale? Zweifellos das europäische Grundgesicht, das unter allem physiognomischen Raffinement ungeändert bleibt. Wir definiren noch zu wenig \Sperr{wesentliche} Unterschiede, wir sind anatomisch noch zwischen engen Grenzen eingeschlossen, vor allem: wir haben noch die gleiche Hautfarbe. Der Europäer ist ein Blassgesicht, darum kommt er über seine „Grundwahrheiten“ nicht hinaus. „Gedanken über kaukasische Vorurtheile“ könnte ein Buch heissen — natürlich wird es nie in Europa geschrieben werden. Höchstens von Einem, der sich von der Meer- und Gebirgssonne lange genug hat bräunen lassen, um uralte, längst verlernte Gefühle in Wiederannäherung zu empfinden: das Für und Wider dunkelgebliebener Rassen den späten, gebleichten gegenüber — tiefes Misstrauen und scheue Anbetung, die im Grunde Schauder vor der Vermischung ist. Ein turanischer Stamm jagt Blondköpfe über die Grenze und opfert weisse Tauben den erzürnten Vorfahren: Verfall in lichten, Gesundheit in dunklen Farben. O ihr gefährlichen, gefährdeten Blassgesichter: ich leide mit euch und mir fehlt kaum eine aus dem Reigen eurer Krankheiten; ich glaube auch an eure „Grundwahrheiten“ und lüge ihnen bessere Gründe unter als ihr. Aber im letzten Innern verberge ich einen guten Willen, zu der früheren, stärkeren Gesundheit anderer Continente und Jahrtausende zurückzureifen, ich hätte nichts dagegen, mit meiner Physiologie wider den Strom zu schwimmen, nämlich ostwärts und zeitaufwärts. Auch die Gelben, auch die Braunen haben ihre „Grundwahrheiten“ \Dots{3}
\par

